\usepackage{PhS_Goodies}% a bunch of newcommands
% packages pour les flotantes
%=============================================================
%
% the packgaes that comes with Latex
%
\usepackage{latexsym}	   % get crazy symbols like Box, Diamond...
\usepackage{ifthen}
\usepackage{makeidx,amsmath, amssymb,  euscript, eufrak}
 
	% beware eufrak MUST be after amssymb
\usepackage{epsfig}   % to feed in figures in .eps mode
\usepackage{array}    % to print nice matrices.
%\usepackage{rotating} % so I can rotate tables
\usepackage{graphics} % graphical packages
%\usepackage{epic}
\usepackage{calc} % in order to calculate values in \newcommand
%
%=============================================================


%=============================================================
%
\vfuzz2pt % Don't report over-full v-boxes if over-edge is small
\hfuzz2pt % Don't report over-full h-boxes if over-edge is small
%
%=============================================================

% numbering paragraphs
\newcounter{parnum}[section] % indexing paragraphs
\newcommand{\Xnewpar}{\vspace{\baselineskip}\stepcounter{parnum}
				\noindent\textbf{\theparnum.}}
%=============================================================
%

\makeindex
%
%=============================================================


%

%=============================================================
%
% the packgaes that were created by author for author himself
%

\usepackage{PhSAbbreviations}
%\usepackage{PhSEquations}

%\usepackage{PhSDefinitions}
%\usepackage{PhSTheorems}
%
%=============================================================



%=============================================================
%
% setting all the counters
% newtheorem implies italic
%

\newtheorem{axiom}{Axiom}[section]
%\newtheorem{definition}{Definition}[section]
%\newtheorem{theorem}{Theorem}[section]
%\newtheorem{lemma}{Lemma}[section]
\newtheorem{proposition}{Proposition}[section]


%
% newcounter keep normal characters
%

\newcounter{examplenum}[section]% indexing examples intra section

\newcommand{\Xexample}[1]{\stepcounter{examplenum}
\noindent \textbf{Example \thesection .\theexamplenum . #1}}
%
%=============================================================



\usepackage{mychapter}
\usepackage{ragged2e}